%----------------------------------------------------------------------------------------
% GONIOMETRISCHE GETALLEN. Alle opgaven zijn eigen werk.
%
% Een hoofdstuk kernvaardigheden (ontwerpregels, deel Kernvaardigheden): beknopte theorie, daarna drill.
% Labels van dit hoofdstuk beginnen met h6-. Alles in graden, zoals in Goniometrie (2de graad).
%----------------------------------------------------------------------------------------

\mychapter{6}{Goniometrische getallen}

\lkr{Herhaling van Driehoeksmeting (3de jaar) en Goniometrie (4de jaar: de goniometrische cirkel en verwante hoeken). Alles in graden; laat de rekenmachine op DEG staan. Bundel 5-8 veronderstellen Goniometrie, hoofdstuk 1. De exacte waarden (bundel 1 en 6) zijn de basis voor de goniometrische functies in de 3de graad.}

\section{Wat je moet kennen}

\begin{theorie}
\naast{\begin{minipage}[t]{8.6cm}
\vspace{0pt}
In een rechthoekige driehoek, voor een scherpe hoek $\alpha$:

\vspace{1mm}
$\formule{\sin\alpha = \dfrac{\text{overstaande}}{\text{schuine zijde}}}$ \quad $\formule{\cos\alpha = \dfrac{\text{aanliggende}}{\text{schuine zijde}}}$

\vspace{1mm}
$\formule{\tan\alpha = \dfrac{\text{overstaande}}{\text{aanliggende}} = \dfrac{\sin\alpha}{\cos\alpha}}$ \qquad $\formule{\sin^2\alpha + \cos^2\alpha = 1}$

\vspace{2mm}
{\renewcommand{\arraystretch}{1.5}\setlength{\tabcolsep}{5pt}%
\begin{tabular}{c|ccccc}
$\alpha$ & $0^\circ$ & $30^\circ$ & $45^\circ$ & $60^\circ$ & $90^\circ$ \\ \hline
$\sin\alpha$ & $0$ & $\tfrac12$ & $\tfrac{\sqrt2}{2}$ & $\tfrac{\sqrt3}{2}$ & $1$ \\
$\cos\alpha$ & $1$ & $\tfrac{\sqrt3}{2}$ & $\tfrac{\sqrt2}{2}$ & $\tfrac12$ & $0$ \\
$\tan\alpha$ & $0$ & $\tfrac{\sqrt3}{3}$ & $1$ & $\sqrt3$ & --
\end{tabular}}
\end{minipage}}{\begin{tikzpicture}[scale=1.55,font=\scriptsize]
  \gonicirkel[t]
  \gonipunt{150}{P}{above left}
  \draw[accent,dashed] (150:1) -- ({cos(150)},0) node[below] {$\cos\alpha$};
  \draw[accent,dashed] (150:1) -- (0,{sin(150)}) node[right] {$\sin\alpha$};
  \draw[grijs,dashed] (0,0) -- (30:1);
  \fill[grijs] (30:1) circle (.03);
  \node[accent] at (.32,.15) {$\alpha$};
\end{tikzpicture}}
\end{theorie}

Op de goniometrische cirkel is $\cos\alpha$ de $x$-coördinaat en $\sin\alpha$ de $y$-coördinaat van het punt $P$ bij $\alpha$. Zo ken je het teken per kwadrant: in het tweede kwadrant is de sinus positief en de cosinus negatief. \textbf{Verwante hoeken} hebben dezelfde getallen op het teken na: $\sin(180^\circ - \alpha) = \sin\alpha$, $\cos(180^\circ - \alpha) = -\cos\alpha$, $\sin(180^\circ + \alpha) = -\sin\alpha$, $\cos(360^\circ - \alpha) = \cos\alpha$.

\textbf{Klassieke fouten.} De rekenmachine in RAD in plaats van DEG. Bij $\sin\alpha = 0{,}3$ alleen $17{,}5^\circ$ geven: ook $180^\circ - 17{,}5^\circ$ heeft die sinus. Het teken vergeten na de grondformule: $\cos\alpha = \pm\sqrt{1 - \sin^2\alpha}$, en het kwadrant beslist.

\voorbeeld{}{Bereken exact $\cos 150^\circ$, en alle hoeken van $0^\circ$ tot $360^\circ$ met $\sin\alpha = 0{,}3$.\label{vb:h6-verwant}}{%
\begin{stappen}
\cos 150^\circ = -\cos 30^\circ = -\frac{\sqrt3}{2} & Supplementaire hoeken: tegengestelde cosinus. \\
\alpha \approx 17{,}5^\circ \text{ of } \alpha \approx 180^\circ - 17{,}5^\circ = 162{,}5^\circ & De rekenmachine geeft er één; de supplementaire hoek is de tweede.
\end{stappen}}

%========================================================================================
\opdrachtenkop

\begin{oefeningen}
\needspace{10\baselineskip}
\oef{h6-o1}

\needspace{10\baselineskip}
\oef{h6-o2}

\needspace{10\baselineskip}
\oef{h6-o3}

\needspace{10\baselineskip}
\oef{h6-o4}

\needspace{10\baselineskip}
\oef{h6-o5}

\needspace{10\baselineskip}
\oef{h6-o6}

\needspace{10\baselineskip}
\oef{h6-o7}

\needspace{10\baselineskip}
\oef{h6-o8}

\needspace{10\baselineskip}
\oef{h6-o9}

\needspace{10\baselineskip}
\oef{h6-o10}
\end{oefeningen}

%----------------------------------------------------------------------------------------
\needspace{14\baselineskip}
\opdrachtsectie{Test jezelf}

\begin{oefeningen}
\oef{h6-t1}

\oef{h6-t2}

\oef{h6-t3}

\oef{h6-t4}

\oef{h6-t5}
\end{oefeningen}

\herhaalschema{%
  \fv{oef:h6-t1}{vb:h6-verwant}{oef:h6-o6}%
  \fzv{oef:h6-t2}{oef:h6-o3}%
  \fzv{oef:h6-t3}{oef:h6-o4}%
  \fzv{oef:h6-t4}{oef:h6-o7}%
  \fv{oef:h6-t5}{vb:h6-verwant}{oef:h6-o8}}

\ifzelfstudie
\section*{Maak je eigen samenvatting}
Sluit dit boek en schrijf in je schrift: sinus, cosinus en tangens in een rechthoekige driehoek, de tabel met de exacte waarden, het teken in elk kwadrant, de verwante hoeken, de grondformule, en hoe je alle hoeken bij een gegeven sinus of cosinus vindt. Vergelijk daarna met het kader.
\fi
